\section{Background and related work}
\label{sec:background}
% ============================================================================

\subsection{Computation in superposition}

A network of hidden width $d$ \emph{represents $F\gg d$ features in superposition} if its
activation $a(x)\in\R^d$ is approximately $\Phi b$ for a feature-encoding matrix
$\Phi\in\R^{d\times F}$ and a sparse Boolean $b$. Following~\cite{hanni2024superposition},
features $f_1,\dots,f_F$ are \emph{$\varepsilon$-linearly represented} by $a$ if there is a
readout $R\in\R^{F\times d}$ with $|r_k\cdot a(x)-f_k(x)|<\varepsilon$ for all $k$ and $x$.

H\"{a}nni et al.\ construct an approximate-linear recursive template for sparse Boolean
circuits, alternating a targeted superpositional AND layer with an error-correction layer, and
certify emulation of circuits of width $\Otilde(d^{3/2})$. Adler and
Shavit~\cite{adler2024superposition} give parameter-description lower bounds and explicit
constructions for computing in superposition, obtaining a near-quadratic capacity by
thresholding after each stage. Adler et al.~\cite{adler2025combinatorial} reframe superposition
combinatorially through feature channel coding, and Kong et al.~\cite{kong2025expand} show
empirically that reducing inter-feature interference at fixed parameter count improves
performance. This paper reproves none of these constructions; it identifies the interface-level
reason the two regimes coexist.

\subsection{Trained toy models of computation in superposition}

A recent empirical line studies whether small trained networks actually compute in
superposition. Braun et al.~\cite{braun2025apd} introduced the compressed-computation task ---
a width-$50$ network asked to compute $100$ sparse ReLU features --- as a testbed. Bhagat et
al.~\cite{bhagat2026cc} argued that the standard $L^2$-trained model largely exploits a mixing
term rather than genuinely computing in superposition, and Ferreira da Silva and
Heimersheim~\cite{ferreiradasilva2026l4} showed that training under an $L^4$ loss elicits a
solution that does compute in superposition, assigning each feature a sparse codeword over
neurons and decoding it with a pseudoinverse of the encoder --- that is, a linear readout
interface in the sense studied here. In parallel,
\cite{ferreiradasilva2026errcorr} report perturbation-based evidence for feature-specific
error correction in large language models, the mechanism whose certified tolerance drives the
$d^{3/2}$ compatibility scale of the H\"{a}nni template. This debate is about which decoding
interface a trained network maintains; the diagnostic proposed here is a direct instrument for
it, and Section~\ref{sec:res-e5} applies it to exactly the models under dispute.

\subsection{Frame theory, Welch bounds and cross-Gramians}
\label{sec:frames}

Classical Welch bounds~\cite{welch1974lower} lower-bound the correlations of overcomplete
unit-norm systems; the maximum-correlation form we use,
$\max_{i\neq j}|\langle\phi_i,\phi_j\rangle|\ge\sqrt{(F-d)/(d(F-1))}$, appears for instance as
Theorem~2.3 of Strohmer and Heath~\cite{strohmer2003grassmannian} in the study of Grassmannian
frames. Waldron~\cite{waldron2003generalized} showed that generalised Welch-bound-equality
sequences are exactly tight frames, and Datta, Howard and Cochran~\cite{datta2012geometry}
derived the entire family of higher-order Welch bounds from a geometric rank/Grammian
argument, of which the rank--trace computation behind Theorem~\ref{thm:floor} is the
first-order instance; see also~\cite{klappenecker2005mub} for the design-theoretic view. For
pairs of sequences, Aceska and Kaczanowski~\cite{aceska2022crossframe} introduced the
cross-frame potential, and Christensen, Datta and Kim~\cite{christensen2020welch} proved
Welch-type bounds for the maximum coherence of dual frame pairs, characterising equality via
equiangular tight frames.

Theorem~\ref{thm:floor} sits in this lineage: it is a unit-diagonal cross-Gramian formulation
assuming only $\operatorname{rank}(M)\le d$ and $M_{ii}=1$ --- in particular, no duality
relation between the two sequences, which, as finite sequences, are of course frames for their
spans. We include its short rank--trace proof only to keep the paper self-contained; the proof
mechanism is standard in this literature, cf.~\cite{datta2012geometry}. The contribution here
is the application to readout interfaces, the gain-normalised form
(Corollary~\ref{cor:calibfree}), and the diagnostic built on top of them.

\subsection{Compressed sensing and sparse recovery}

The threshold-recovery results used below are standard coherence-based and random-dictionary
support-recovery arguments. Their role is to make explicit a distinction relevant to neural
computation: small-error linear readout requires uniformly small coordinate error, whereas
threshold recovery only requires the aggregate interference from the active support to remain
below a margin.

The collision frontier of Section~\ref{sec:frontier} sits in this neighbourhood and we place it
carefully, because two established literatures ask adjacent questions.

The first is group testing. Recovering a sparse Boolean $b$ from $\Phi b$ is exactly
\emph{quantitative} group testing when $\Phi$ is a real measurement matrix, and the combinatorial
theory of superimposed, disjunct and separable codes descends from Kautz and
Singleton~\cite{kautz1964superimposed}. That theory asks how many tests suffice to identify the
\emph{whole} support, by \emph{any} decoder, usually for a matrix to be designed. We ask something
weaker and more specific: for a matrix we are given, and one coordinate of it, is that coordinate's
presence decidable by an \emph{affine} function of $\Phi b$? The gap is real in both directions.
A matrix can be identifiable in the group-testing sense while $\kappa_i\le s$, because
$\kappa_i\le s$ only requires a collision between \emph{convex combinations} of admissible states,
not an exact coincidence of two integer ones. Conversely $\kappa_i>s$ certifies a linear probe
with a margin, which identifiability alone does not supply. The restriction to affine decoders is
not a weakness of the analysis but the object of interest, since linear probing is how features
are read in practice.

The second is the null space property. Writing $z'$ for $z$ extended by zero at coordinate $i$,
the constraint $\Phi_{-i}z=\phi_i$ says precisely that $z'-e_i\in\ker\Phi$, so $\kappa_i$ is a
statement about kernel vectors normalised at one coordinate --- the same geometry as the null
space property and as the neighbourliness of randomly projected
polytopes~\cite{donoho2005neighborliness}. Two differences matter. The null space property is
quantified over all supports and certifies recovery by $\ell_1$ minimisation, a nonlinear decoder;
$\kappa_i$ is per-coordinate and certifies an affine one. And the null space property is hard to
verify, which is why the literature tests it by semidefinite
relaxation~\cite{daspremont2011nullspace}, whereas $\kappa_i$ is the exact value of one linear
programme --- a consequence of asking the weaker, one-coordinate question rather than of a better
algorithm. The box $\norm{z}_\infty\le1/\alpha$, which comes from bounded activations and has no
counterpart in the null space property, turns out to bind rarely, so for most features the
frontier is a minimum-$\ell_1$-representation quantity and the tools of that literature apply
directly.

\subsection{Sparse autoencoders}

Sparse autoencoder work scales dictionary learning to millions of
latents~\cite{cunningham2023sparse,gao2024scaling,lieberum2024gemma} and studies
reconstruction--sparsity trade-offs~\cite{michaud2025manifolds}. Such results motivate the
importance of overcomplete internal representations, but dictionary size is a reconstructive
quantity and does not measure the recursive Boolean interface capacities studied here. This
paper answers one of the open questions listed by Sharkey et al.~\cite{sharkey2025open}
concerning what theoretical insight can be gleaned from computation performed natively in
superposition.

\paragraph{Delineation from a prior sparse-autoencoder study.} The closest prior work measures
sparse autoencoders with an overlapping toolset and a different question.
\citet{borobia2026pruning}, in \emph{Neurocomputing}, train TopK autoencoders on weight-pruned
language models and ask which features survive compression, using reconstruction fidelity and
feature-level survival statistics as the instrument. That paper is the reference point for the
dictionary measurements here, and the division of labour between the two is worth stating plainly
because the objects overlap and the quantities do not. This paper trains no autoencoder, prunes
nothing, and measures no reconstruction. It
takes decoder matrices as given --- including ones that paper never examined, such as the
randomly-initialised control of Section~\ref{sec:dictionaries} --- and computes a different quantity:
the cross-talk a linear reader of that decoder must accept, its exact decomposition into leverage
dispersion, and the sparsity at which affine recovery must fail. No model, dictionary, measurement or
result is shared between the two, and neither paper's conclusions depend on the other's.

% ============================================================================
